% 贯穿图邻接矩阵的平方：(A^2)_{1,4}=a_{1,2}a_{2,4}+a_{1,3}a_{3,4}=2，对应行走1,2,4与1,3,4。
\begin{tikzpicture}[>=Stealth,
  v/.style={circle,draw=FigureStroke,fill=FigureBlueFill,line width=.8pt,minimum size=6.5mm,inner sep=0pt},
  e/.style={draw=FigureGrid,line width=.8pt},
  wb/.style={draw=FigureBlue,line width=1.5pt,->,shorten >=1pt,shorten <=1pt},
  wg/.style={draw=FigureGreen,line width=1.5pt,->,shorten >=1pt,shorten <=1pt},
  cell/.style={draw=FigureGrid,line width=.5pt,minimum width=6mm,minimum height=5.2mm,inner sep=0pt},
  lab/.style={text=FigureMuted,minimum width=6mm,minimum height=5.2mm,inner sep=0pt},
  hl/.style={cell,fill=FigureYellowFill,draw=FigureYellow,line width=.9pt},
  note/.style={font=\FigureNoteFont\fontsize{7.5}{9}\selectfont,text=FigureMuted,align=center}]
\node[v] (1) at (-1.3,0) {$1$}; \node[v] (2) at (0,.95) {$2$};
\node[v] (3) at (0,-.95) {$3$}; \node[v] (4) at (1.3,0) {$4$};
\draw[e] (2)--(3);
\draw[wb] (1)--(2); \draw[wb] (2)--(4);
\draw[wg] (1)--(3); \draw[wg] (3)--(4);
\node[text=FigureBlue] at (-.95,.75) {$1,2,4$};
\node[text=FigureGreen] at (-.95,-.75) {$1,3,4$};
\node at (0,-1.8) {(a) 从$1$到$4$长度为$2$的行走};
\matrix[matrix of math nodes,nodes={cell},column sep=-.5pt,row sep=-.5pt] (M) at (4.6,.1) {
  |[lab,draw=none]| & |[lab]|1 & |[lab]|2 & |[lab]|3 & |[lab]|4 \\
  |[lab]|1 & 2 & 1 & 1 & |[hl]|2 \\
  |[lab]|2 & 1 & 3 & 2 & 1 \\
  |[lab]|3 & 1 & 2 & 3 & 1 \\
  |[lab]|4 & 2 & 1 & 1 & 2 \\};
\node at (4.6,-1.8) {(b) $\bm A^2$};
\node[anchor=west] at (6.4,.55) {$(\bm A^2)_{1,4}=\sum_{\ell}a_{1,\ell}\,a_{\ell,4}$};
\node[anchor=west] at (6.4,-.05) {$=\textcolor{FigureBlue}{a_{1,2}a_{2,4}}+\textcolor{FigureGreen}{a_{1,3}a_{3,4}}$};
\node[anchor=west] at (6.4,-.65) {$=1+1=2$};
\end{tikzpicture}
